\subsection{Other designs}

Some instruments replace the array detector with a mechanical scan, which then requires an angle calibration.\citep{gonzalez2020development,zainurin2023developing,hosker2017demonstrating} A scanner with a DVD grating covered 400--700~nm in steps of 1.52~nm; its absorbances differed from those of a Cary~60 by 0--12.84\% for coumarin and by 0.52--4.14\% in a mercury-ion assay.\citep{gonzalez2020development} Brick-built and optical-breadboard instruments let students see scanning and alignment at work,\citep{bougotrobin2016optimization,davis2018open} although stray light lowered the calibration slope of the unenclosed breadboard.\citep{davis2018open}

Other designs try out new approaches to fabrication or detection. In a 500--700~nm instrument with printed lenses and a prism, a filter with a nominal width of 10~nm appeared about 18~nm wide, a measure of how much the assembly broadened the band.\citep{berglund2022enabling} A UV instrument converted transmitted light at 205~nm into scintillator fluorescence that a camera could pick up; its sensitivity to nitrate was 27\% lower than that of the reference with a bandpass filter, and 85\% lower without one.\citep{booysen2021water}

\subsection{Choosing an architecture}

Choosing an architecture starts from the spectral information the measurement needs and from the geometry of the sample. Fixed bands are sufficient for well-defined assays, whereas resolved spectra reveal spectral structure and unexpected absorbers. Custom gratings give full control of the optics, integrated modules save alignment work, and multispectral chips offer compact measurements at selected bands. Which of these is appropriate comes down to whether it can meet the measurement requirement with the construction and calibration resources available.
\FloatBarrier
